\section{Setting, Estimators, and Efficiency}
\label{sec:setting}

\paragraph{Sampling model.}
We observe $n$ labeled units $(Z_i,Y_i)$ and $N$ unlabeled units $\tilde Z_j$, where $Z=(f(X),\,\cdot)$ collects the information available to the analyst (the prediction alone, or the prediction together with covariates $X$). All units are i.i.d. from one population, and labeled and unlabeled samples are independent. Let $a=n/(n+N)$ denote the labeled fraction and $\theta=\mathbb E[Y]$ the target (a mean, including prevalence for binary $Y$). Write $m(Z)=\mathbb E[Y\mid Z]$, $\sigma^2(Z)=\mathrm{Var}(Y\mid Z)$, and let $L(Z)$ be the best linear predictor of $Y$ from $Z$ (with intercept).

\paragraph{A common form.}
For any calibrator $\hat m$ (a function fitted on labeled data), define
\begin{equation}
\hat\theta_{\hat m}
=\bar Y_L+(1-a)\Big(\bar{\hat m}_U-\bar{\hat m}_L\Big)
=\underbrace{\bar Y_L-(1-a)\bar{\hat m}_L}_{\text{labeled}}
+\underbrace{(1-a)\bar{\hat m}_U}_{\text{unlabeled}}\,,
\label{eq:common}
\end{equation}
where $\bar{\hat m}_L$ and $\bar{\hat m}_U$ are sample means over labeled and unlabeled units. Equivalently, $\hat\theta$ is the average of $\hat m$ over the pooled sample plus the labeled residual correction $\bar Y_L-\bar{\hat m}_L$. Taking $\hat m=f$ and $\hat m$ linear gives PPI and PPI++ (with $\lambda=(1-a)\beta$, $\beta=\mathrm{Cov}(Y,f)/\mathrm{Var}(f)$, which is the PPI++ value $\lambda^\star=\mathrm{Cov}(Y,f)/\{\mathrm{Var}(f)(1+n/N)\}$).

\paragraph{Cross-fitted calibrator.}
Split the labeled units into $K$ folds. Fit $\hat m_{-k}$ on the labeled units outside fold $k$. For labeled unit $i$ in fold $k(i)$ set $\hat m_L(Z_i)=\hat m_{-k(i)}(Z_i)$ (out-of-fold), and for unlabeled units set $\hat m_U(\tilde Z_j)=K^{-1}\sum_k\hat m_{-k}(\tilde Z_j)$. The estimator is \eqref{eq:common} with these values, and the variance estimate is
\begin{equation}
\widehat V=\frac{\widehat{\mathrm{Var}}_L\{Y-(1-a)\hat m_L\}}{n}+\frac{(1-a)^2\,\widehat{\mathrm{Var}}_U(\hat m_U)}{N}\,,
\qquad \mathrm{CI}=\hat\theta\pm z_{1-\alpha/2}\sqrt{\widehat V}.
\label{eq:cfvar}
\end{equation}
With $\hat m=L$ fitted on all labeled data, \eqref{eq:cfvar} reduces to the PPI++ variance formula.

\subsection{Three facts about this estimator class}

\begin{proposition}[Exact unbiasedness under cross-fitting]
\label{prop:unbiased}
If the fold fits are exchangeable (e.g., random folds and a deterministic learner) and integrable, then $\mathbb E\,\hat\theta=\theta$ for every $n$ and $N$, whatever the learner is and however large its error.
\end{proposition}
\begin{proof}[Proof sketch]
$\mathbb E\bar Y_L=\theta$. Each out-of-fold value has the same mean $\mu_{\hat m}=\mathbb E\hat m_{-k}(Z)$ by exchangeability, so $\mathbb E\bar{\hat m}_L=\mathbb E\bar{\hat m}_U=\mu_{\hat m}$, and the correction term has mean zero.
\end{proof}

Unbiasedness does not depend on the learner, which is why cross-fitting matters here. Without it, $\hat m_L$ and $Y_L$ share data and the labeled residual term carries an in-sample bias.

\begin{proposition}[Asymptotic variance]
\label{prop:var}
Suppose the fold fits converge in $L_2$ to a fixed function $\tilde m$ (possibly $\neq m$), with finite fourth moments of $Y$ and $\tilde m$. Then $\sqrt n(\hat\theta-\theta)\Rightarrow\mathcal N(0,V(\tilde m))$, where
\begin{equation}
n\,V(\tilde m)=\mathbb E[\sigma^2(Z)]+a\,\mathrm{Var}\{m(Z)\}+(1-a)\,\mathrm{Var}\{\tilde m(Z)-m(Z)\}.
\label{eq:Vcal}
\end{equation}
\end{proposition}
\begin{proof}[Proof sketch]
Write $\delta=\tilde m-m$. Linearizing \eqref{eq:common} gives $\hat\theta-\theta\approx(\bar Y_L-\theta)-(1-a)(\bar{\tilde m}_L-\mu)+(1-a)(\bar{\tilde m}_U-\mu)$. The labeled and unlabeled pieces are independent, so $V=\mathrm{Var}\{Y-(1-a)\tilde m\}/n+(1-a)^2\mathrm{Var}(\tilde m)/N$. Since $Y-(1-a)\tilde m=e+am-(1-a)\delta$ with $e=Y-m$ uncorrelated with any function of $Z$, expanding and using $(1-a)^2/N=a(1-a)/n$, the coefficient of $\mathrm{Var}(m)$ is $a^2+a(1-a)=a$, that of $\mathrm{Cov}(m,\delta)$ is $-2a(1-a)+2a(1-a)=0$, and that of $\mathrm{Var}(\delta)$ is $(1-a)^2+a(1-a)=1-a$, which gives \eqref{eq:Vcal}.
\end{proof}

Two consequences matter. First, the calibrator's \emph{error} enters the variance only through $\mathrm{Var}(\tilde m-m)$, with no bias term. Second, the penalty carries the factor $1-a=N/(n+N)$, so it vanishes when $N=0$.

\begin{theorem}[Efficiency bound and the PPI++ gap]
\label{thm:bound}
Let $\mathcal M_Z$ be the model in which $Z$ is the only information used, $(Z,Y)$ is otherwise unrestricted, and labels are missing completely at random with probability $a$. The semiparametric efficiency bound for $\theta$ is
\begin{equation}
V^\star(Z)=\frac{\mathbb E[\sigma^2(Z)]+a\,\mathrm{Var}\{m(Z)\}}{n},
\label{eq:bound}
\end{equation}
attained by \eqref{eq:common} with $\tilde m=m$. Let $\Delta_L=\mathbb E\{m(Z)-L(Z)\}^2$. Then the PPI++ estimator (linear $\hat m$ in $Z$) has
\begin{equation}
V_{\mathrm{PPI}\text{++}}-V^\star(Z)=\frac{(1-a)\,\Delta_L}{n}.
\label{eq:gap}
\end{equation}
\end{theorem}
\begin{proof}[Proof sketch]
For the observation $O=(R,RY,Z)$ with $R\sim\mathrm{Bernoulli}(a)$ independent of $(Z,Y)$, the efficient influence function under missing-at-random outcomes is $\psi=\frac{R}{a}\{Y-m(Z)\}+m(Z)-\theta$ (Robins, Rotnitzky and Zhao, 1994; Tsiatis, 2006). Its variance per unit is $\mathbb E\sigma^2/a+\mathrm{Var}(m)$, and with $n+N$ units and $a(n+N)=n$ this gives \eqref{eq:bound}. For PPI++, $\tilde m=L$, so $\delta=L-m$ has mean zero and $\mathrm{Var}(\delta)=\Delta_L$; apply Proposition~\ref{prop:var}.
\end{proof}

\begin{corollary}[When can the interval shrink?]
\label{cor:when}
Write $\rho=V^\star/V_{\mathrm{PPI}\text{++}}$. The asymptotic CI width ratio of the efficient interval to PPI++ is $\sqrt\rho$, where
\[
1-\rho=\frac{(1-a)\,\Delta_L}{\mathbb E(Y-L)^2+a\,\mathrm{Var}(L)}.
\]
Hence: (i) PPI++ is efficient for the information $Z$ iff $m(Z)$ is affine in $Z$ ($\Delta_L=0$); (ii) if $N=0$ there is no gain; (iii) as $a\to0$ the gain tends to $\Delta_L/\mathbb E(Y-L)^2$, the share of the linear calibrator's residual variance explained by the nonlinear part of $m$.
\end{corollary}

\begin{corollary}[Covariates beyond the prediction]
\label{cor:cov}
Let $Z^+=(f,X)$ and $Z=f$. Then $V^\star(Z)-V^\star(Z^+)=\frac{(1-a)}{n}\mathbb E\{m(Z^+)-m(Z)\}^2$. Using covariates helps by exactly the same type of term as nonlinear calibration.
\end{corollary}

\begin{corollary}[Learned calibrators]
\label{cor:learned}
Let $\hat m$ converge to $\tilde m$. The cross-fitted estimator beats PPI++ asymptotically iff
\[
\mathrm{Var}(\tilde m-m)<\Delta_L,
\]
i.e., the calibrator's error must be smaller than the linear calibrator's misfit. The asymptotic width ratio is $\sqrt{1+(1-a)\{\mathrm{Var}(\tilde m-m)-\Delta_L\}/(nV_{\mathrm{PPI}\text{++}})}$.
\end{corollary}

These statements are asymptotic and conditional on the chosen information set. They do not say that a valid interval can never be shorter than PPI++ under other assumptions, for example bounded outcomes or exact small-sample constructions. The claim is about regular Wald-type intervals built from the same $(Z_L,Y_L,Z_U)$.
